\textbf{Phần 1. Trắc nghiệm lựa chọn}
\setcounterpageref{ex}{0}

\begin{ex} %[5D12N70]
Tỉ số của hai số $4$ và $9$ là
\choice
{\True $\dfrac{4}{9}$}
{$\dfrac{9}{4}$}
{$4 \times 9$}
{$36$}
\loigiai{
Tỉ số là $\dfrac{4}{9}$.\\
\textbf{Đáp án của câu hỏi là A}
}
\end{ex}

\begin{ex} %[5D12N70]
Số thập phân $0{,}35$ viết thành tỉ số phần trăm là
\choice
{$3{,}5\%$}
{\True $35\%$}
{$350\%$}
{$0{,}35\%$}
\loigiai{
$0{,}35 = 35\%$.\\
\textbf{Đáp án của câu hỏi là B}
}
\end{ex}

\begin{ex} %[5D12H70]
Tìm $20\%$ của $150\text{ kg}$:
\choice
{$20\text{ kg}$}
{$40\text{ kg}$}
{\True $30\text{ kg}$}
{$50\text{ kg}$}
\loigiai{
$150 \times 20 : 100 = 30\text{ (kg)}$.\\
\textbf{Đáp án của câu hỏi là C}
}
\end{ex}

\begin{ex} %[5D12H70]
Một lớp học có $40$ học sinh, trong đó có $24$ bạn nữ. Tỉ số phần trăm học sinh nam so với cả lớp là
\choice
{$60\%$}
{\True $40\%$}
{$50\%$}
{$45\%$}
\loigiai{
Số bạn nam: $40 - 24 = 16$ (bạn).\\
Tỉ số phần trăm: $16 : 40 = 40\%$.\\
\textbf{Đáp án của câu hỏi là B}
}
\end{ex}

\begin{ex} %[5D12H70]
Tổng của hai số là $120$, tỉ số của hai số là $\dfrac{1}{3}$. Số bé là
\choice
{$40$}
{$60$}
{$90$}
{\True $30$}
\loigiai{
$120 : (1 + 3) \times 1 = 30$.\\
\textbf{Đáp án của câu hỏi là D}
}
\end{ex}

\begin{ex} %[5D12H70]
Hiệu của hai số là $45$, biết số lớn gấp $4$ lần số bé. Số lớn là
\choice
{$15$}
{\True $60$}
{$45$}
{$75$}
\loigiai{
Số lớn: $45 : (4 - 1) \times 4 = 60$.\\
\textbf{Đáp án của câu hỏi là B}
}
\end{ex}

\begin{ex} %[5D12V70]
Trên bản đồ tỉ lệ $1 : 20,000$, quãng đường từ nhà An đến trường đo được $6\text{ cm}$. Khoảng cách thực tế là
\choice
{$120\text{ m}$}
{$12\text{ km}$}
{$120\text{ km}$}
{\True $1{,}2\text{ km}$}
\loigiai{
$6 \times 20,000 = 120,000\text{ cm} = 1{,}2\text{ km}$.\\
\textbf{Đáp án của câu hỏi là D}
}
\end{ex}

\begin{ex} %[5D12V70]
Một chiếc xe máy có giá $25\,000\,000$ đồng, nhân dịp khuyến mãi được giảm giá $8\%$. Số tiền khách hàng được giảm là
\choice
{$1\,500\,000$ đồng}
{$2\,500\,000$ đồng}
{\True $2\,000\,000$ đồng}
{$1\,800\,000$ đồng}
\loigiai{
$25\,000\,000 \times 8 : 100 = 2\,000\,000$ (đồng).\\
\textbf{Đáp án của câu hỏi là C}
}
\end{ex}

\begin{ex} %[5D12V70]
Một mảnh đất hình chữ nhật có chu vi $100\text{ m}$, chiều rộng bằng $\dfrac{2}{3}$ chiều dài. Diện tích mảnh đất đó là
\choice
{$500\text{ m}^2$}
{$700\text{ m}^2$}
{$800\text{ m}^2$}
{\True $600\text{ m}^2$}
\loigiai{
Nửa chu vi: $50\text{ m}$.\\
Chiều dài: $50 : (2 + 3) \times 3 = 30\text{ m}$.\\
Chiều rộng: $50 - 30 = 20\text{ m}$.\\
Diện tích: $30 \times 20 = 600\text{ (m}^2\text{)}$.\\
\textbf{Đáp án của câu hỏi là D}
}
\end{ex}

\begin{ex} %[5D12V70]
Một người gửi tiết kiệm $60\,000\,000$ đồng với lãi suất $6\%$ một năm. Sau một năm người đó nhận được số tiền lãi là
\choice
{\True $3\,600\,000$ đồng}
{$3\,000\,000$ đồng}
{$4\,000\,000$ đồng}
{$6\,000\,000$ đồng}
\loigiai{
$60\,000\,000 \times 6 : 100 = 3\,600\,000$ (đồng).\\
\textbf{Đáp án của câu hỏi là A}
}
\end{ex}

\begin{ex} %[5D12V70]
Biết $45\%$ của một số là $180$. Số đó là
\choice
{$81$}
{$360$}
{\True $400$}
{$450$}
\loigiai{
$180 : 45 \times 100 = 400$.\\
\textbf{Đáp án của câu hỏi là C}
}
\end{ex}

\begin{ex} %[5D12C70]
Một cửa hàng bán một chiếc đồng hồ với giá $600\,000$ đồng thì được lãi $20\%$ so với giá vốn. Giá vốn của chiếc đồng hồ đó là
\choice
{\True $500\,000$ đồng}
{$480\,000$ đồng}
{$520\,000$ đồng}
{$550\,000$ đồng}
\loigiai{
Giá bán chiếm $100\% + 20\% = 120\%$ giá vốn.\\
Giá vốn: $600\,000 : 120 \times 100 = 500\,000$ (đồng).\\
\textbf{Đáp án của câu hỏi là A}
}
\end{ex}

\textbf{Phần 2. Trắc nghiệm đúng sai}
\setcounterpageref{ex}{0}

\begin{ex} %[5D12H70]
Một nông trại có $300$ con vịt và $200$ con gà. Xét tính đúng sai:
\choiceTF
{\True Tỉ số giữa số vịt và số gà là $\dfrac{3}{2}$.}
{\True Số vịt chiếm $60\%$ tổng đàn gia cầm.}
{Số gà chiếm $30\%$ tổng đàn gia cầm.}
{\True Số gà chiếm $40\%$ tổng đàn gia cầm.}
\loigiai{
\begin{itemize}
    \item Tỉ số vịt và gà: $300 : 200 = \dfrac{3}{2}$. Đúng.
    \item Tổng đàn: $500$ con $\Rightarrow$ Vịt: $300 : 500 = 60\%$. Đúng.
    \item Gà: $200 : 500 = 40\% \neq 30\%$. Ý 3 Sai, Ý 4 Đúng.
\end{itemize}
\begin{itemize}
    \item a) \textbf{Đúng}.
    \item b) \textbf{Đúng}.
    \item c) \textbf{Sai}.
    \item d) \textbf{Đúng}.
\end{itemize}
}
\end{ex}

\begin{ex} %[5D12V70]
Hai ngăn sách có tất cả $180$ cuốn, ngăn thứ nhất bằng $\dfrac{4}{5}$ ngăn thứ hai. Xét tính đúng sai:
\choiceTF
{\True Tổng số phần bằng nhau là $9$ phần.}
{\True Ngăn thứ nhất có $80$ cuốn sách.}
{Ngăn thứ hai có $90$ cuốn sách.}
{\True Ngăn thứ hai nhiều hơn ngăn thứ nhất $20$ cuốn sách.}
\loigiai{
\begin{itemize}
    \item Tổng số phần: $4 + 5 = 9$. Đúng.
    \item Ngăn 1: $180 : 9 \times 4 = 80$ cuốn. Đúng.
    \item Ngăn 2: $180 - 80 = 100$ cuốn $\neq 90$ cuốn. Sai.
    \item Hiệu: $100 - 80 = 20$ cuốn. Đúng.
\end{itemize}
\begin{itemize}
    \item a) \textbf{Đúng}.
    \item b) \textbf{Đúng}.
    \item c) \textbf{Sai}.
    \item d) \textbf{Đúng}.
\end{itemize}
}
\end{ex}

\textbf{Phần 3. Trả lời ngắn}
\setcounterpageref{ex}{0}

\begin{ex} %[5D12H70]
Tìm tỉ số phần trăm của $18$ và $40$. (Chỉ ghi số).
\shortans{45}
\loigiai{
\textbf{Đáp số:} $45$\\
$18 : 40 = 45\%$.
}
\end{ex}

\begin{ex} %[5D12H70]
Tìm $35\%$ của $240$.
\shortans{84}
\loigiai{
\textbf{Đáp số:} $84$\\
$240 \times 35 : 100 = 84$.
}
\end{ex}

\begin{ex} %[5D12V70]
Tổng hai số là $144$, tỉ số là $\dfrac{3}{5}$. Tìm số lớn.
\shortans{90}
\loigiai{
\textbf{Đáp số:} $90$\\
$144 : (3 + 5) \times 5 = 90$.
}
\end{ex}

\begin{ex} %[5D12V70]
Hiệu hai số là $56$, tỉ số là $\dfrac{2}{9}$. Tìm số bé.
\shortans{16}
\loigiai{
\textbf{Đáp số:} $16$\\
$56 : (9 - 2) \times 2 = 16$.
}
\end{ex}

\begin{ex} %[5D12V70]
Trên bản đồ tỉ lệ $1 : 5,000$, một thửa ruộng hình chữ nhật có chiều dài $4\text{ cm}$. Tính chiều dài thật theo mét.
\shortans{200}
\loigiai{
\textbf{Đáp số:} $200$\\
$4 \times 5,000 = 20,000\text{ cm} = 200\text{ m}$.
}
\end{ex}

\begin{ex} %[5D12C70]
Một sản phẩm có giá gốc là $400\,000$ đồng. Người bán nâng giá lên $20\%$, sau đó lại giảm giá $10\%$ trên giá mới. Hỏi giá bán cuối cùng là bao nhiêu nghìn đồng?
\shortans{432}
\loigiai{
\textbf{Đáp số:} $432$\\
Giá sau khi nâng: $400 \times 1{,}2 = 480$ nghìn đồng.\\
Giá sau khi giảm 10\%: $480 \times 0{,}9 = 432$ nghìn đồng.
}
\end{ex}
